\documentclass[11pt,a4paper]{article}
\usepackage{amsmath, amssymb, amsfonts}
\usepackage{geometry}
\usepackage{graphicx}
\usepackage{algorithm}
\usepackage{algpseudocode}
\usepackage{hyperref}
\usepackage{booktabs}
\usepackage{caption}
\usepackage[backend=biber,style=ieee,sorting=none]{biblatex}

\geometry{margin=1in}

% Embedded references
\usepackage{filecontents}
\begin{filecontents}{\jobname.bib}
@article{BlackScholes1973,
  author={F. Black and M. Scholes},
  title={The Pricing of Options and Corporate Liabilities},
  journal={The Journal of Political Economy},
  volume={81},
  number={3},
  pages={637--654},
  year={1973}
}
@article{Merton1976,
  author={R. C. Merton},
  title={Option pricing when underlying stock returns are discontinuous},
  journal={Journal of Financial Economics},
  volume={3},
  number={1-2},
  pages={125--144},
  year={1976}
}
@book{Oksendal2003,
  author={B. {\O}ksendal},
  title={Stochastic Differential Equations: An Introduction with Applications},
  publisher={Springer},
  year={2003}
}
@book{Klebaner2005,
  author={F. C. Klebaner},
  title={Introduction to Stochastic Calculus with Applications},
  publisher={Imperial College Press},
  year={2005}
}
@article{Bian2016,
  author={X. Bian and C. Kim and G. E. Karniadakis},
  title={111 years of Brownian motion},
  journal={Soft Matter},
  volume={12},
  number={30},
  pages={6331--6346},
  year={2016}
}
@book{Ross2012,
  author={S. M. Ross},
  title={Simulations},
  publisher={Academic Press},
  year={2012}
}
@misc{Hyndman2024,
  author={C. Hyndman},
  title={Continuous-time Finance},
  note={PowerPoint presented at MACF 402, Mathematical and Computational Finance II, Concordia University, 2024}
}
@article{Artzner1999,
  author={P. Artzner and F. Delbaen and J.-M. Eber and D. Heath},
  title={Coherent measures of risk},
  journal={Mathematical Finance},
  volume={9},
  number={3},
  pages={203--228},
  year={1999}
}
\end{filecontents}

\addbibresource{\jobname.bib}

\title{5-day Value-at-Risk Estimation for a European Call Option \\ Using Delta-Normal Approximation and Importance Sampling}
\author{Zakaria El Yassini}
\date{December 28, 2024}

\begin{document}

\maketitle

\begin{abstract}
This work presents a rigorous framework for estimating the \textit{Value-at-Risk} (VaR) of a long European call option using the \textit{Delta-Normal Approximation}. The underlying asset is modeled as a geometric Brownian motion on a filtered probability space $(\Omega, \mathcal{F}, \{\mathcal{F}_t\}_{t \ge 0}, \mathbb{Q})$ under the risk-neutral measure $\mathbb{Q}$. The approach combines Black-Scholes option pricing \cite{BlackScholes1973}, Delta computation \cite{Hyndman2024}, and Importance Sampling using a Girsanov change of measure to assess tail risk. The 5-day VaR is formally defined following Artzner et al. \cite{Artzner1999}.
\end{abstract}

\section{Introduction}
Value-at-Risk (VaR) quantifies the minimal loss a portfolio may incur over a given time horizon at a specified confidence level \cite{Artzner1999}. Let $(\Omega, \mathcal{F}, \{\mathcal{F}_t\}, \mathbb{P})$ be a filtered probability space supporting a standard Brownian motion. Then, for portfolio loss $L$ over horizon $\Delta t$, the $\alpha$-level VaR is defined as:

\begin{equation*}
\text{VaR}_\alpha = 
\left\{
\begin{array}{l}
\inf\{x \in \mathbb{R} : \mathbb{P}(L > x) \le \alpha\}, \\[1mm]
L = V_0 - V_{\Delta t}
\end{array}
\right.
\end{equation*}

This work focuses on a long European call option, estimating the 5-day VaR via the Delta-Normal approximation. Importance Sampling is applied to better capture tail events \cite{Ross2012,Hyndman2024}.

\section{Theoretical Framework}

\subsection{Asset Dynamics}
The underlying risky asset $S_t^1$ evolves according to the Black-Scholes model \cite{BlackScholes1973,Merton1976}:

\begin{equation*}
\left\{
\begin{array}{l}
dS_t^1 =\mu S_t^1 \, dt + \sigma S_t^1 \, dB_t, \\[1mm]
S_0^1 > 0
\end{array}
\right.
\end{equation*}
where $\mu$ is the drift, $\sigma$ is volatility, and $B_t$ is a standard Brownian motion adapted to the filtration $\{\mathcal{F}_t\}$. The risk-free bank account evolves as:

\begin{equation*}
\left\{
\begin{array}{l}
dS_t^0 = r S_t^0 dt \\[1mm]
S_0^0=1.
\end{array}
\right.
\end{equation*}

\subsection{European Call Option Pricing}
Let $C(t,F)$ denote the price of a European call option with strike $K$ and maturity $T$. It satisfies the Black-Scholes PDE:

\begin{equation*}
\begin{cases}
\displaystyle \frac{\partial C}{\partial t} + \frac{1}{2} \sigma^2 F^2 \frac{\partial^2 C}{\partial F^2} - r C = 0, & 0 \le t < T, \\[3mm]
C(T, F(T,T)) = \max(F(T,T) - K, 0)
\end{cases}
\end{equation*}
The analytical solution is:

\begin{equation*}
C(t,S_t) = S_t N(d_1) - K e^{-r(T-t)} N(d_2),
\end{equation*}
with
\begin{equation*}
\begin{cases}
d_1 = \displaystyle \frac{\ln(S_t/K) + \left(r + \frac{1}{2}\sigma^2\right)(T-t)}{\sigma \sqrt{T-t}},\\[2mm]
d_2 = d_1 - \sigma \sqrt{T-t}.
\end{cases}
\end{equation*}


\subsection{Delta and Delta-Normal Approximation for VaR}

The \textit{Delta} of an option quantifies the sensitivity of the option's value to small changes in the underlying asset price. For a European call option, the Delta is given by

\begin{equation*}
\Delta = \frac{\partial C}{\partial S_0^1} 
= e^{-rT} \, \mathbb{E}^\mathbb{Q} \left[ \mathbf{1}_{\{S_T^1>K\}} \frac{S_T^1}{S_0^1} \right]
\text{\hfill \cite{Hyndman2024}}
\end{equation*}
where the expectation is taken under the risk-neutral measure $\mathbb{Q}$. This expression captures the probabilistic impact of the terminal payoff on the option's current value, reflecting the likelihood that the option finishes in-the-money. \\
\\
Under the \textit{Delta-Normal Approximation}, for sufficiently small movements in the underlying asset, the change in the option value $\Delta V$ can be linearized as:
\begin{equation*}
\begin{cases}
\Delta V \approx \Delta \cdot V_0(S_0) \cdot \Delta S,\\[2mm]
\Delta S \sim N(0, \sigma^2 \Delta t).
\end{cases}
\end{equation*}
where $\Delta S$ represents the change in the underlying asset price over the risk horizon $\Delta t$, assumed to be normally distributed. This simplification allows us to model the portfolio loss using a univariate normal distribution, making the calculation of Value-at-Risk computationally efficient.\\
\\
Consequently, the $\alpha$-level Value-at-Risk over the horizon $\Delta t$ is computed as

\begin{equation*}
\begin{cases}
\text{VaR}_\alpha = - z_\alpha \, \Delta \, \sigma \sqrt{\Delta t},\\[2mm]
z_\alpha = N^{-1}(1-\alpha).
\end{cases}
\text{\hfill \cite{Artzner1999}}
\end{equation*}
where $z_\alpha$ is the $\alpha$-quantile of the standard normal distribution. In this context, VaR represents the **minimum loss that will not be exceeded with probability $1-\alpha$**, providing a coherent measure of tail risk for risk management purposes.


\subsection{Importance Sampling via Girsanov}
To capture tail events more effectively, we shift the drift under a new measure $\mathbb{Q}^{\text{IS}}$ \cite{Hyndman2024}:

\begin{equation}
\mu_{\text{IS}} = r - 0.5\sigma^2 + \text{shift factor} \cdot \sigma.
\end{equation}
The Radon-Nikodym derivative gives the path weights:

\begin{equation}
w_i = \frac{d\mathbb{Q}}{d\mathbb{Q}^{\text{IS}}} = \exp(-( \mu_{\text{IS}} - r)\Delta t).
\end{equation}

\section{Algorithms}

\begin{algorithm}[H]
\caption{Delta-Normal VaR Estimation}
\begin{algorithmic}[1]
\State \textbf{Input:} $S_0,K,r,\sigma,T,\Delta t,\alpha$
\State Compute option Delta $\Delta$
\State Compute $z_\alpha = \text{qnorm}(1-\alpha)$
\State Compute VaR: $\text{VaR} = - z_\alpha \Delta \sigma \sqrt{\Delta t}$
\State \textbf{Output:} VaR
\end{algorithmic}
\end{algorithm}

\begin{algorithm}[H]
\caption{Importance Sampling for VaR Breaks}
\begin{algorithmic}[1]
\State \textbf{Input:} $S_0,K,r,\sigma,T,\Delta t,\text{shift factor},N$
\State Initialize $Breaks = 0$
\For{$i = 1$ to $N$}
    \State Simulate $S_i^{\text{IS}} = S_0 \exp((\mu_{\text{IS}}-0.5\sigma^2)\Delta t + \sigma \sqrt{\Delta t} Z_i)$
    \State Compute option value $V_i$
    \State Compute weight $w_i = \exp(-( \mu_{\text{IS}} - r)\Delta t)$
    \If{$V_i - V_0 < \text{VaR}$}
        \State $Breaks \gets Breaks + w_i$
    \EndIf
\EndFor
\State Estimate $\hat{\alpha}_{\text{IS}} = Breaks / N$
\State \textbf{Output:} $\hat{\alpha}_{\text{IS}}$
\end{algorithmic}
\end{algorithm}

\section{Numerical Example}
\begin{table}[h!]
\centering
\begin{tabular}{lc}
\toprule
Parameter & Value \\
\midrule
$S_0$ & 100 \\
$K$ & 100 \\
$r$ & 0.01 \\
$\sigma$ & 0.5 \\
$T$ & 1 year \\
$\Delta t$ & 5/365 \\
$\alpha$ & 0.05 \\
$N$ & $10^6$ \\
\bottomrule
\end{tabular}
\caption{Simulation parameters}
\end{table}

\subsection{Results}
The 5-day Value-at-Risk (VaR) for a long European call option, calculated using the Delta-Normal Approximation, is:

\[
\text{VaR}_{5\text{-day}} = -5.837263.
\]
This indicates that, under the model assumptions, the minimum expected loss in the worst 5\% of cases over a 5-day horizon is approximately \$5.84. In other words, the portfolio is expected to lose at least this amount with 5\% probability \cite{Artzner1999}.\\  
By incorporating Importance Sampling to better capture the tail of the distribution of asset returns, the estimated probability of observing a loss exceeding this VaR threshold is:

\[
\hat{\alpha}_{\text{IS}} = 0.02256787.
\]
This implies that in the simulated scenarios, only 2.26\% of outcomes exceeded the computed VaR, which is lower than the nominal 5\% level, reflecting the effectiveness of tail-focused sampling in capturing extreme events. These results highlight that using Importance Sampling can provide a more accurate estimate of tail risk, and help risk managers make informed decisions about capital allocation and potential hedging strategies.\\
\\
Overall, the combination of the Delta-Normal approximation with Importance Sampling demonstrates a pragmatic balance between computational efficiency and the need for realistic tail-risk assessment, providing actionable insights for risk management in options portfolios.


\section{Discussion and Improvements}
\begin{itemize}
    \item \textbf{Real-Time Data Integration:} Implementing live market data sources (e.g., via \texttt{yfinance}) allows continuous updating of VaR estimates to reflect current market conditions.
    \item \textbf{Model Enhancements:} The Delta-Normal Approximation assumes normally distributed returns, which underestimates extreme events. Incorporating stochastic volatility models (e.g., GARCH), fat-tailed distributions, or Monte Carlo simulations can provide more robust risk estimates.
    \item \textbf{Dynamic Risk Management:} VaR estimates can feed into dynamic hedging strategies, portfolio rebalancing, and risk limit frameworks to reduce the likelihood of large, unexpected losses.
\end{itemize}

\section{Conclusion}
The Delta-Normal Approximation offers a computationally efficient method to estimate VaR for European call options, capturing the primary sensitivity of the option to underlying price movements. Incorporating Importance Sampling via a Girsanov change of measure significantly improves the assessment of tail risk, providing a more accurate estimate of the probability and magnitude of extreme losses. Future work can enhance this framework by integrating real-time data, modeling non-normal return dynamics, and implementing fully dynamic risk management strategies to ensure robust portfolio protection under adverse market conditions.
\newpage
\printbibliography

\end{document}

